\documentclass[12pt,a4paper]{article}
\usepackage[UTF8]{ctex}
\usepackage{amsmath,amssymb}
\usepackage{geometry}
\usepackage{listings}
\usepackage{xcolor}
\usepackage{tcolorbox}
\usepackage{enumitem}
\usepackage{hyperref}

\geometry{margin=2.5cm}

% 定义颜色
\definecolor{lightgray}{RGB}{240,240,240}
\definecolor{myblue}{RGB}{70,130,180}
\definecolor{mygreen}{RGB}{34,139,34}
\definecolor{myred}{RGB}{220,20,60}
\definecolor{myorange}{RGB}{255,140,0}
\definecolor{myyellow}{RGB}{218,165,32}
\definecolor{mypurple}{RGB}{128,0,128}

% C++ 代码高亮设置
\lstset{
	language=C++,
	basicstyle=\ttfamily\small,
	keywordstyle=\color{myblue}\bfseries,
	commentstyle=\color{gray},
	stringstyle=\color{teal},
	numbers=left,
	numberstyle=\tiny\color{gray},
	stepnumber=1,
	numbersep=5pt,
	backgroundcolor=\color{lightgray!30},
	frame=single,
	breaklines=true,
	breakatwhitespace=false,
	tabsize=4,
	showspaces=false,
	showstringspaces=false,
	showtabs=false,
	captionpos=b,
	morekeywords={long long, vector, pair, unordered\_map}
}

% 自定义盒子环境 - 修正版
\newtcolorbox{bluebox}{
	colback=myblue!10!white,
	colframe=myblue!70!black,
	fonttitle=\bfseries,
	arc=4pt
}

\newtcolorbox{greenbox}{
	colback=mygreen!10!white,
	colframe=mygreen!70!black,
	fonttitle=\bfseries,
	arc=4pt
}

\newtcolorbox{redbox}{
	colback=myred!10!white,
	colframe=myred!70!black,
	fonttitle=\bfseries,
	arc=4pt
}

\newtcolorbox{orangebox}{
	colback=myorange!10!white,
	colframe=myorange!70!black,
	fonttitle=\bfseries,
	arc=4pt
}

\newtcolorbox{yellowbox}{
	colback=myyellow!10!white,
	colframe=myyellow!70!black,
	fonttitle=\bfseries,
	arc=4pt
}

\title{\Huge \textbf{二分查找算法完全指南}}
\author{算法学习笔记}
\date{\today}

\begin{document}
	
	\maketitle
	
	\tableofcontents
	\newpage
	
	\section{二分查找的核心概念}
	
	\subsection{什么是二分查找？}
	
	二分查找（Binary Search）是一种在有序序列中快速查找特定元素的算法。其核心思想是"减而治之"——每次都将搜索范围缩小一半。
	
	\begin{bluebox}
		\textbf{形象理解}
		
		想象你在翻一本词典找某个单词：
		\begin{enumerate}
			\item 打开到中间位置
			\item 如果目标单词应该在前面，就把后半部分丢掉
			\item 如果目标单词应该在后面，就把前半部分丢掉
			\item 重复这个过程，直到找到目标
		\end{enumerate}
	\end{bluebox}
	
	\subsection{为什么时间复杂度是$O(\log n)$？}
	
	每次查找都能排除一半的数据：
	\begin{itemize}
		\item 第1次：剩 $n/2$ 个数据
		\item 第2次：剩 $n/4$ 个数据
		\item 第3次：剩 $n/8$ 个数据
		\item ...
		\item 第$k$次：剩 $n/2^k$ 个数据
	\end{itemize}
	
	当 $n/2^k = 1$ 时，$k = \log_2 n$，因此最多需要 $\log_2 n$ 次比较。
	
	\subsection{二分查找的关键——"二段性"}
	
	\textbf{重要理解：}二分查找的本质不是数据"有序"，而是数据具有"二段性"——我们可以用一个条件把数据分成"满足条件"和"不满足条件"的两段。
	
	\begin{yellowbox}
		\textbf{二段性示例}
		
		\begin{itemize}
			\item 在有序数组中，以某个值为界，分为"小于目标值"和"大于等于目标值"
			\item 在旋转数组中，以旋转点为界，分为"较大的有序段"和"较小的有序段"
			\item 在峰值问题中，分为"在上升坡上"和"在下降坡上"
		\end{itemize}
	\end{yellowbox}
	
	\section{整数二分查找的两种经典模板}
	
	\subsection{模板一：查找第一个满足条件的元素（左边界）}
	
	\textbf{适用场景：}找到第一个满足条件的元素位置
	
	\begin{greenbox}
		\textbf{模板一核心记忆}
		
		\textbf{"左边界四步法":}
		\begin{enumerate}
			\item 初始化：$l = 0, r = n - 1$
			\item 循环条件：$while(l < r)$
			\item 中点计算：$mid = (l + r) / 2$（向下取整）
			\item 判断更新：
			\begin{itemize}
				\item 如果 $mid$ 满足条件：$r = mid$（不排除mid，向左收缩）
				\item 如果 $mid$ 不满足条件：$l = mid + 1$（排除mid，向右走）
			\end{itemize}
		\end{enumerate}
	\end{greenbox}
	
	\begin{lstlisting}[caption={模板1：查找第一个 >= x 的元素}]
		#include <iostream>
		#include <vector>
		using namespace std;
		
		// 在有序数组arr中，查找第一个 >= x 的元素下标
		// 如果不存在，返回 -1
		int findFirstGreaterOrEqual(vector<int>& arr, int x) {
			int l = 0, r = arr.size() - 1;
			
			while (l < r) {
				int mid = l + (r - l) / 2;  // 等价于 (l + r) / 2，但防止溢出
				
				if (arr[mid] >= x) {
					// mid满足条件（>=x），可能是答案，往左找更小的
					r = mid;
				} else {
					// mid不满足条件（<x），mid肯定不是答案
					l = mid + 1;
				}
			}
			
			// while结束时 l == r，这个位置就是答案
			if (arr[l] >= x) return l;
			return -1;  // 所有元素都 < x
		}
		
		// 使用示例
		int main() {
			vector<int> arr = {1, 3, 3, 3, 5, 7, 9};
			int x = 3;
			int pos = findFirstGreaterOrEqual(arr, x);
			cout << "第一个 >= " << x << " 的位置: " << pos << endl;
			// 输出: 第一个 >= 3 的位置: 1
			return 0;
		}
	\end{lstlisting}
	
	\subsection{模板二：查找最后一个满足条件的元素（右边界）}
	
	\textbf{适用场景：}找到最后一个满足条件的元素位置
	
	\begin{orangebox}
		\textbf{模板二核心记忆}
		
		\textbf{"右边界四步法":}
		\begin{enumerate}
			\item 初始化：$l = 0, r = n - 1$
			\item 循环条件：$while(l < r)$
			\item 中点计算：$mid = (l + r + 1) / 2$（向上取整，防止死循环）
			\item 判断更新：
			\begin{itemize}
				\item 如果 $mid$ 满足条件：$l = mid$（不排除mid，向右收缩）
				\item 如果 $mid$ 不满足条件：$r = mid - 1$（排除mid，向左走）
			\end{itemize}
		\end{enumerate}
	\end{orangebox}
	
	\begin{lstlisting}[caption={模板2：查找最后一个 <= x 的元素}]
		#include <iostream>
		#include <vector>
		using namespace std;
		
		// 在有序数组arr中，查找最后一个 <= x 的元素下标
		// 如果不存在，返回 -1
		int findLastLessOrEqual(vector<int>& arr, int x) {
			int l = 0, r = arr.size() - 1;
			
			while (l < r) {
				int mid = l + (r - l + 1) / 2;  // 注意：+1是关键，防止死循环
				
				if (arr[mid] <= x) {
					// mid满足条件（<=x），可能是答案，往右找更大的
					l = mid;
				} else {
					// mid不满足条件（>x），mid肯定不是答案
					r = mid - 1;
				}
			}
			
			if (arr[l] <= x) return l;
			return -1;  // 所有元素都 > x
		}
		
		// 使用示例
		int main() {
			vector<int> arr = {1, 3, 3, 3, 5, 7, 9};
			int x = 3;
			int pos = findLastLessOrEqual(arr, x);
			cout << "最后一个 <= " << x << " 的位置: " << pos << endl;
			// 输出: 最后一个 <= 3 的位置: 3
			return 0;
		}
	\end{lstlisting}
	
	\subsection{为什么模板二需要 $mid = (l + r + 1) / 2$？}
	
	这是为了防止死循环。考虑 $l = 0, r = 1$ 的情况：
	
	\begin{redbox}
		\textbf{死循环分析}
		
		\begin{itemize}
			\item \textbf{如果使用} $mid = (l + r) / 2$：
			\begin{itemize}
				\item $mid = (0 + 1) / 2 = 0$
				\item 如果满足条件，执行 $l = mid = 0$
				\item 区间还是 $[0, 1]$，没变化！进入死循环
			\end{itemize}
			\item \textbf{如果使用} $mid = (l + r + 1) / 2$：
			\begin{itemize}
				\item $mid = (0 + 1 + 1) / 2 = 1$
				\item 如果满足条件，执行 $l = mid = 1$
				\item 区间变成 $[1, 1]$，退出循环
			\end{itemize}
		\end{itemize}
	\end{redbox}
	
	\section{浮点数二分查找}
	
	浮点数二分不需要考虑整数边界问题，只需要控制精度。
	
	\begin{lstlisting}[caption={浮点数二分：求平方根}]
		#include <iostream>
		#include <cmath>
		using namespace std;
		
		// 计算x的平方根，精确到1e-6
		double sqrtBinary(double x) {
			double l = 0, r = max(1.0, x);  // 注意处理x<1的情况
			
			// 方法1：控制绝对误差
			while (r - l > 1e-8) {  // 精度要求
				double mid = (l + r) / 2;
				if (mid * mid >= x) {
					r = mid;
				} else {
					l = mid;
				}
			}
			return l;
		}
		
		// 方法2：固定循环次数（更常用）
		double sqrtBinary2(double x) {
			double l = 0, r = max(1.0, x);
			
			for (int i = 0; i < 100; i++) {  // 循环100次足够精确
				double mid = (l + r) / 2;
				if (mid * mid >= x) {
					r = mid;
				} else {
					l = mid;
				}
			}
			return l;
		}
		
		int main() {
			double x = 2.0;
			cout << "sqrt(2) = " << sqrtBinary(x) << endl;
			cout << "实际值: " << sqrt(2) << endl;
			// 输出: sqrt(2) ≈ 1.414214
			return 0;
		}
	\end{lstlisting}
	
	\section{STL中的二分查找函数}
	
	C++ STL提供了现成的二分查找函数：
	
	\begin{lstlisting}[caption={STL二分查找函数使用示例}]
		#include <iostream>
		#include <algorithm>
		#include <vector>
		using namespace std;
		
		int main() {
			vector<int> arr = {1, 3, 3, 3, 5, 7, 9};
			
			// 1. binary_search：判断元素是否存在
			bool exists = binary_search(arr.begin(), arr.end(), 3);
			cout << "3是否存在: " << (exists ? "是" : "否") << endl;
			
			// 2. lower_bound：查找第一个 >= x 的位置（左边界模板）
			auto left = lower_bound(arr.begin(), arr.end(), 3);
			cout << "第一个 >= 3 的位置: " << (left - arr.begin()) << endl;
			cout << "对应的值: " << *left << endl;
			
			// 3. upper_bound：查找第一个 > x 的位置
			auto right = upper_bound(arr.begin(), arr.end(), 3);
			cout << "第一个 > 3 的位置: " << (right - arr.begin()) << endl;
			cout << "对应的值: " << *right << endl;
			
			// 4. equal_range：同时获取左右边界
			auto range = equal_range(arr.begin(), arr.end(), 3);
			cout << "3的范围: [" << (range.first - arr.begin()) 
			<< ", " << (range.second - arr.begin()) << ")" << endl;
			
			// 总结：如果我们要找3出现的次数
			int count = right - left;
			cout << "3出现的次数: " << count << endl;
			
			return 0;
		}
	\end{lstlisting}
	
	\section{深入理解：二段性与check函数}
	
	\subsection{什么是好的check函数？}
	
	二分查找的核心是设计一个\textbf{check函数}，使得数据具有"二段性"：
	
	\begin{bluebox}
		\textbf{check函数设计原则}
		
		\begin{itemize}
			\item \textbf{单调性：}前面一段不满足，后面一段满足（或反之）
			\item \textbf{明确性：}对于任意位置，能明确判断是否满足条件
			\item \textbf{目标性：}check函数的临界点就是我们想找的答案
		\end{itemize}
	\end{bluebox}
	
	\subsection{如何构造check函数？}
	
	\begin{lstlisting}[caption={check函数的设计模式}]
		// 模式1：找第一个满足条件的（左边界）
		// check函数：判断是否满足目标性质
		bool check(vector<int>& arr, int mid, int target) {
			return arr[mid] >= target;  // 满足条件的性质
		}
		
		int binarySearchLeft(vector<int>& arr, int target) {
			int l = 0, r = arr.size() - 1;
			while (l < r) {
				int mid = l + (r - l) / 2;
				if (check(arr, mid, target)) {
					r = mid;  // 满足条件，往左找
				} else {
					l = mid + 1;  // 不满足，往右找
				}
			}
			return l;
		}
		
		// 模式2：找最后一个满足条件的（右边界）
		int binarySearchRight(vector<int>& arr, int target) {
			int l = 0, r = arr.size() - 1;
			while (l < r) {
				int mid = l + (r - l + 1) / 2;
				if (check(arr, mid, target)) {
					l = mid;  // 满足条件，往右找
				} else {
					r = mid - 1;  // 不满足，往左找
				}
			}
			return l;
		}
	\end{lstlisting}
	
	\section{经典例题详解}
	
	\subsection{例题1：在排序数组中查找元素的第一个和最后一个位置}
	
	\textbf{题目描述：}给定一个升序排列的整数数组和一个目标值，找出目标值在数组中的开始位置和结束位置。
	
	\begin{lstlisting}[caption={例题1：查找元素的起始和结束位置}]
		#include <iostream>
		#include <vector>
		using namespace std;
		
		vector<int> searchRange(vector<int>& nums, int target) {
			if (nums.empty()) return {-1, -1};
			
			// 第一次二分：找第一个 >= target 的位置（左边界）
			int l = 0, r = nums.size() - 1;
			while (l < r) {
				int mid = l + (r - l) / 2;
				if (nums[mid] >= target) {
					r = mid;
				} else {
					l = mid + 1;
				}
			}
			
			// 检查是否找到了target
			if (nums[l] != target) return {-1, -1};
			int left_bound = l;
			
			// 第二次二分：找最后一个 <= target 的位置（右边界）
			r = nums.size() - 1;  // l不变，只重置r
			while (l < r) {
				int mid = l + (r - l + 1) / 2;  // 注意+1
				if (nums[mid] <= target) {
					l = mid;
				} else {
					r = mid - 1;
				}
			}
			
			return {left_bound, l};
		}
		
		int main() {
			vector<int> nums = {5, 7, 7, 8, 8, 10};
			int target = 8;
			vector<int> result = searchRange(nums, target);
			
			cout << "数组: ";
			for (int num : nums) cout << num << " ";
			cout << endl;
			cout << "查找目标: " << target << endl;
			cout << "范围: [" << result[0] << ", " << result[1] << "]" << endl;
			// 输出: [3, 4]
			
			return 0;
		}
	\end{lstlisting}
	
	\subsection{例题2：寻找峰值}
	
	\textbf{题目描述：}峰值元素是指其值大于左右相邻值的元素。给定一个数组，返回任意一个峰值元素的索引。
	
	\textbf{解题思路：}虽然数组无序，但具有二段性：
	\begin{itemize}
		\item 如果 $nums[mid] < nums[mid+1]$，那么mid在上升坡上，峰值一定在右侧
		\item 否则，mid在下降坡或峰顶，峰值在左侧（包括mid）
	\end{itemize}
	
	\begin{lstlisting}[caption={例题2：寻找峰值}]
		#include <iostream>
		#include <vector>
		using namespace std;
		
		int findPeakElement(vector<int>& nums) {
			int l = 0, r = nums.size() - 1;
			
			while (l < r) {
				int mid = l + (r - l) / 2;
				
				// 比较mid和mid+1，判断在上升坡还是下降坡
				if (nums[mid] > nums[mid + 1]) {
					// mid在下降坡，峰值在左侧（包括mid）
					r = mid;
				} else {
					// mid在上升坡，峰值在右侧（不包括mid）
					l = mid + 1;
				}
			}
			
			return l;  // l就是峰值位置
		}
		
		int main() {
			vector<int> nums = {1, 2, 1, 3, 5, 6, 4};
			int peak = findPeakElement(nums);
			
			cout << "数组: ";
			for (int num : nums) cout << num << " ";
			cout << endl;
			cout << "峰值位置: " << peak << ", 峰值大小: " << nums[peak] << endl;
			// 可能的输出: 峰值位置: 5, 峰值大小: 6
			
			return 0;
		}
	\end{lstlisting}
	
	\subsection{例题3：搜索旋转排序数组}
	
	\textbf{题目描述：}一个原本升序的数组在某个点进行了旋转，搜索目标值的位置。
	
	例如：$[4,5,6,7,0,1,2]$ 在原数组$[0,1,2,4,5,6,7]$的基础上旋转了4个位置。
	
	\textbf{解题思路：}二段性体现在——总能找到一段是有序的，判断target是否在这段中。
	
	\begin{lstlisting}[caption={例题3：搜索旋转排序数组}]
		#include <iostream>
		#include <vector>
		using namespace std;
		
		int searchRotated(vector<int>& nums, int target) {
			int l = 0, r = nums.size() - 1;
			
			while (l <= r) {  // 注意：这里是 <=
				int mid = l + (r - l) / 2;
				
				if (nums[mid] == target) {
					return mid;  // 找到了
				}
				
				// 判断哪边是有序的
				if (nums[l] <= nums[mid]) {
					// 左半部分有序
					if (nums[l] <= target && target < nums[mid]) {
						// target在左半部分
						r = mid - 1;
					} else {
						// target在右半部分
						l = mid + 1;
					}
				} else {
					// 右半部分有序
					if (nums[mid] < target && target <= nums[r]) {
						// target在右半部分
						l = mid + 1;
					} else {
						// target在左半部分
						r = mid - 1;
					}
				}
			}
			
			return -1;  // 没找到
		}
		
		int main() {
			vector<int> nums = {4, 5, 6, 7, 0, 1, 2};
			int target = 0;
			
			cout << "旋转数组: ";
			for (int num : nums) cout << num << " ";
			cout << endl;
			cout << "查找目标: " << target << endl;
			cout << "位置: " << searchRotated(nums, target) << endl;
			// 输出: 位置: 4
			
			return 0;
		}
	\end{lstlisting}
	
	\subsection{例题4：找第一个错误版本}
	
	\textbf{题目描述：}有$n$个版本$[1, n]$，已知在某个版本之后都是错误的。用最少的调用次数找到第一个错误的版本。
	
	\begin{lstlisting}[caption={例题4：第一个错误版本}]
		#include <iostream>
		using namespace std;
		
		// 模拟的API：判断版本是否正确
		bool isBadVersion(int version) {
			// 假设从版本4开始都是错误的
			return version >= 4;
		}
		
		int firstBadVersion(int n) {
			int l = 1, r = n;
			
			while (l < r) {
				int mid = l + (r - l) / 2;
				
				if (isBadVersion(mid)) {
					// 当前版本是错误的，第一个错误版本可能是mid或更早
					r = mid;
				} else {
					// 当前版本正确，第一个错误版本在后面
					l = mid + 1;
				}
			}
			
			return l;  // l就是第一个错误版本
		}
		
		int main() {
			int n = 10;
			int bad = firstBadVersion(n);
			
			cout << "总版本数: " << n << endl;
			cout << "第一个错误版本: " << bad << endl;
			// 输出: 第一个错误版本: 4
			
			return 0;
		}
	\end{lstlisting}
	
	\subsection{例题5：计算x的平方根}
	
	\textbf{题目描述：}实现整数平方根函数，计算并返回$x$的平方根，只保留整数部分。
	
	\begin{lstlisting}[caption={例题5：x的平方根（整数版本）}]
		#include <iostream>
		using namespace std;
		
		int mySqrt(int x) {
			if (x == 0) return 0;
			
			int l = 1, r = x;
			
			while (l < r) {
				int mid = l + (r - l + 1) / 2;  // 使用右边界模板
				
				// 判断mid的平方是否 <= x
				if (mid <= x / mid) {  // 用除法避免溢出
					// mid^2 <= x，满足条件，可能还有更大的答案
					l = mid;
				} else {
					// mid^2 > x，不满足，答案在左边
					r = mid - 1;
				}
			}
			
			return l;
		}
		
		int main() {
			int x1 = 8, x2 = 16, x3 = 2147395599;
			
			cout << "sqrt(" << x1 << ") = " << mySqrt(x1) << endl;  // 2
			cout << "sqrt(" << x2 << ") = " << mySqrt(x2) << endl;  // 4
			cout << "sqrt(" << x3 << ") = " << mySqrt(x3) << endl;  // 46339
			
			return 0;
		}
	\end{lstlisting}
	
	\subsection{例题6：在D天内完成所有任务}
	
	\textbf{题目描述：}传送带上有$n$个包裹，第$i$个包裹重量为$weights[i]$。要在$D$天内按顺序发送所有包裹，求传送带的最小运载能力。
	
	\textbf{解题思路：}二分答案——能力值越大，所需天数越少。找到第一个满足天数$\leq D$的能力。这就是典型的"对答案进行二分"。
	
	\begin{lstlisting}[caption={例题6：传送带运载能力}]
		#include <iostream>
		#include <vector>
		#include <algorithm>
		using namespace std;
		
		// 计算以capacity的运载能力需要多少天
		int daysNeeded(vector<int>& weights, int capacity) {
			int days = 1;  // 至少需要1天
			int current_sum = 0;
			
			for (int weight : weights) {
				if (current_sum + weight > capacity) {
					// 当前包裹放不下了，需要新的一天
					days++;
					current_sum = weight;
				} else {
					current_sum += weight;
				}
			}
			
			return days;
		}
		
		int shipWithinDays(vector<int>& weights, int D) {
			// 运载能力的范围：至少是最大包裹重量，最大是所有包裹总重量
			int left = *max_element(weights.begin(), weights.end());
			int right = 0;
			for (int w : weights) right += w;
			
			// 二分查找最小运载能力（左边界模板）
			while (left < right) {
				int mid = left + (right - left) / 2;
				
				if (daysNeeded(weights, mid) <= D) {
					// 以mid的运力能在D天内完成，说明mid可能太大了
					// 尝试更小的运力
					right = mid;
				} else {
					// 以mid的运力不能在D天内完成，必须增大运力
					left = mid + 1;
				}
			}
			
			return left;
		}
		
		int main() {
			vector<int> weights = {1, 2, 3, 4, 5, 6, 7, 8, 9, 10};
			int D = 5;
			
			int capacity = shipWithinDays(weights, D);
			cout << "包裹重量: ";
			for (int w : weights) cout << w << " ";
			cout << endl;
			cout << "需要在 " << D << " 天内完成" << endl;
			cout << "最小运载能力: " << capacity << endl;
			// 输出: 15
			
			return 0;
		}
	\end{lstlisting}
	
	\section{难题进阶——有深度的二分问题}
	
	\subsection{进阶题1：寻找两个正序数组的中位数}
	
	\textbf{题目描述：}给定两个大小分别为 $m$ 和 $n$ 的正序数组 $nums1$ 和 $nums2$。找出这两个正序数组的中位数，要求时间复杂度为 $O(\log(m+n))$。
	
	\textbf{解题思路：}这道题是二分的经典难题。核心思想是：
	\begin{itemize}
		\item 在较短的数组上二分查找"分割线"的位置
		\item 分割线左边元素个数 = 分割线右边元素个数（或左边多1个）
		\item 满足交叉大小关系：$nums1[i-1] \leq nums2[j]$ 且 $nums2[j-1] \leq nums1[i]$
	\end{itemize}
	
	\begin{lstlisting}[caption={进阶题1：两个有序数组的中位数}]
		#include <iostream>
		#include <vector>
		#include <algorithm>
		#include <climits>
		using namespace std;
		
		double findMedianSortedArrays(vector<int>& nums1, vector<int>& nums2) {
			// 确保nums1是较短的数组，减少二分次数
			if (nums1.size() > nums2.size()) {
				return findMedianSortedArrays(nums2, nums1);
			}
			
			int m = nums1.size(), n = nums2.size();
			int totalLeft = (m + n + 1) / 2;  // 左半部分应该有多少个元素
			
			// 在nums1上二分查找合适的分割线
			int l = 0, r = m;
			
			while (l < r) {
				int i = l + (r - l + 1) / 2;  // nums1的分割线位置
				int j = totalLeft - i;        // nums2的分割线位置
				
				if (nums1[i - 1] > nums2[j]) {
					// nums1分割线左边太大了，需要左移
					r = i - 1;
				} else {
					// 满足条件，尝试右移找更大的i
					l = i;
				}
			}
			
			int i = l;  // nums1的分割位置
			int j = totalLeft - i;  // nums2的分割位置
			
			// 处理边界情况
			int nums1LeftMax = (i == 0) ? INT_MIN : nums1[i - 1];
			int nums1RightMin = (i == m) ? INT_MAX : nums1[i];
			int nums2LeftMax = (j == 0) ? INT_MIN : nums2[j - 1];
			int nums2RightMin = (j == n) ? INT_MAX : nums2[j];
			
			// 计算中位数
			if ((m + n) % 2 == 1) {
				// 奇数个元素，中位数是左半部分最大值
				return max(nums1LeftMax, nums2LeftMax);
			} else {
				// 偶数个元素，中位数是左右部分最大最小值的平均
				double left = max(nums1LeftMax, nums2LeftMax);
				double right = min(nums1RightMin, nums2RightMin);
				return (left + right) / 2.0;
			}
		}
		
		int main() {
			vector<int> nums1 = {1, 3};
			vector<int> nums2 = {2};
			
			cout << "数组1: 1 3" << endl;
			cout << "数组2: 2" << endl;
			cout << "中位数: " << findMedianSortedArrays(nums1, nums2) << endl;
			// 输出: 2.0
			
			vector<int> nums3 = {1, 2};
			vector<int> nums4 = {3, 4};
			cout << "\n数组1: 1 2" << endl;
			cout << "数组2: 3 4" << endl;
			cout << "中位数: " << findMedianSortedArrays(nums3, nums4) << endl;
			// 输出: 2.5
			
			return 0;
		}
	\end{lstlisting}
	
	\subsection{进阶题2：分割数组的最大值}
	
	\textbf{题目描述：}给定一个非负整数数组和一个整数 $m$，将数组分成 $m$ 个非空的连续子数组。设计一个算法使得这 $m$ 个子数组各自和的最大值最小。
	
	\textbf{解题思路：}
	\begin{itemize}
		\item 二分答案：子数组和的最大值在 $[\max(nums), \sum(nums)]$ 之间
		\item check函数：判断能否在不超过该最大和的情况下分成 $\leq m$ 个子数组
		\item 找左边界：第一个能分成 $\leq m$ 组的最大和
	\end{itemize}
	
	\begin{lstlisting}[caption={进阶题2：分割数组的最大值最小化}]
		#include <iostream>
		#include <vector>
		#include <algorithm>
		#include <numeric>
		using namespace std;
		
		// 判断在最大和为maxSum的情况下，能否分成 <= m 个子数组
		bool canSplit(vector<int>& nums, int m, int maxSum) {
			int count = 1;  // 当前分成的子数组个数
			int currentSum = 0;
			
			for (int num : nums) {
				if (currentSum + num > maxSum) {
					// 当前子数组加不下这个数了，需要新的一组
					count++;
					currentSum = num;
					if (count > m) return false;
				} else {
					currentSum += num;
				}
			}
			
			return true;
		}
		
		int splitArray(vector<int>& nums, int m) {
			// 二分范围：至少是最大元素，至多是所有元素之和
			long long left = *max_element(nums.begin(), nums.end());
			long long right = accumulate(nums.begin(), nums.end(), 0LL);
			
			// 找左边界：第一个能分成 <= m 组的最大和
			while (left < right) {
				long long mid = left + (right - left) / 2;
				
				if (canSplit(nums, m, mid)) {
					// mid可行，尝试更小的最大值
					right = mid;
				} else {
					// mid太小了，分成的组数 > m，需要增大
					left = mid + 1;
				}
			}
			
			return left;
		}
		
		int main() {
			vector<int> nums = {7, 2, 5, 10, 8};
			int m = 2;
			
			cout << "数组: 7 2 5 10 8" << endl;
			cout << "分割成 " << m << " 份" << endl;
			cout << "最小的最大子数组和: " << splitArray(nums, m) << endl;
			// 输出: 18 (分为 [7,2,5] 和 [10,8])
			
			return 0;
		}
	\end{lstlisting}
	
	\subsection{进阶题3：找出第k小的距离对}
	
	\textbf{题目描述：}给定一个整数数组，返回所有数对之间的绝对差值中第 $k$ 小的值。
	
	\textbf{解题思路：}
	\begin{itemize}
		\item 首先排序数组
		\item 二分答案：差值在 $[0, \max-\min]$ 之间
		\item check函数：计算差值 $\leq mid$ 的数对个数（用双指针或二分）
		\item 找第一个满足"差值 $\leq mid$ 的数对个数 $\geq k$"的值
	\end{itemize}
	
	\begin{lstlisting}[caption={进阶题3：第k小的距离对}]
		#include <iostream>
		#include <vector>
		#include <algorithm>
		using namespace std;
		
		// 计算差值 <= mid 的数对个数
		int countPairs(vector<int>& nums, int mid) {
			int count = 0;
			int n = nums.size();
			
			// 双指针法，对于每个左指针i，找到最远的右指针j
			for (int i = 0, j = 0; i < n; i++) {
				while (j < n && nums[j] - nums[i] <= mid) {
					j++;
				}
				// 对于固定的i，有j-i-1个数满足条件
				count += j - i - 1;
			}
			
			return count;
		}
		
		int smallestDistancePair(vector<int>& nums, int k) {
			// 先排序
			sort(nums.begin(), nums.end());
			
			int n = nums.size();
			int left = 0, right = nums[n - 1] - nums[0];
			
			// 二分查找第一个满足 countPairs(mid) >= k 的值
			while (left < right) {
				int mid = left + (right - left) / 2;
				
				if (countPairs(nums, mid) >= k) {
					// 差值 <= mid 的数对已经 >= k，说明第k小的差值 <= mid
					right = mid;
				} else {
					// 差值 <= mid 的数对不够k个，第k小的差值 > mid
					left = mid + 1;
				}
			}
			
			return left;
		}
		
		int main() {
			vector<int> nums = {1, 3, 1};
			int k = 1;
			
			cout << "数组: 1 3 1" << endl;
			cout << "第 " << k << " 小的距离对: " 
			<< smallestDistancePair(nums, k) << endl;
			// 所有距离对: (1,3)=2, (1,1)=0, (3,1)=2
			// 排序后: 0, 2, 2, 第1小是0
			
			nums = {1, 1, 1};
			k = 2;
			cout << "\n数组: 1 1 1" << endl;
			cout << "第 " << k << " 小的距离对: " 
			<< smallestDistancePair(nums, k) << endl;
			// 所有距离对都是0，第2小是0
			
			nums = {1, 6, 1};
			k = 3;
			cout << "\n数组: 1 6 1" << endl;
			cout << "第 " << k << " 小的距离对: " 
			<< smallestDistancePair(nums, k) << endl;
			// 距离对: (1,6)=5, (1,1)=0, (6,1)=5
			// 排序后: 0, 5, 5, 第3小是5
			
			return 0;
		}
	\end{lstlisting}
	
	\subsection{进阶题4：在排序矩阵中查找第k小的元素}
	
	\textbf{题目描述：}给定一个 $n \times n$ 矩阵，每行和每列元素均按升序排列。找到矩阵中第 $k$ 小的元素。
	
	\textbf{解题思路：}
	\begin{itemize}
		\item 二分答案：第k小元素在 $[matrix[0][0], matrix[n-1][n-1]]$ 之间
		\item check函数：计算矩阵中 $\leq mid$ 的元素个数
		\item 利用矩阵的有序性，在$O(n)$时间内完成计数
	\end{itemize}
	
	\begin{lstlisting}[caption={进阶题4：排序矩阵中的第k小元素}]
		#include <iostream>
		#include <vector>
		using namespace std;
		
		// 计算矩阵中 <= mid 的元素个数
		int countLessEqual(vector<vector<int>>& matrix, int mid) {
			int n = matrix.size();
			int count = 0;
			
			// 从左下角开始统计
			int row = n - 1, col = 0;
			
			while (row >= 0 && col < n) {
				if (matrix[row][col] <= mid) {
					// 当前列的所有上方元素都 <= mid
					count += row + 1;
					col++;  // 向右移动
				} else {
					row--;  // 向上移动
				}
			}
			
			return count;
		}
		
		int kthSmallest(vector<vector<int>>& matrix, int k) {
			int n = matrix.size();
			int left = matrix[0][0];
			int right = matrix[n - 1][n - 1];
			
			// 找左边界：第一个满足 countLessEqual(mid) >= k 的值
			while (left < right) {
				int mid = left + (right - left) / 2;
				
				if (countLessEqual(matrix, mid) >= k) {
					right = mid;
				} else {
					left = mid + 1;
				}
			}
			
			return left;
		}
		
		int main() {
			vector<vector<int>> matrix = {
				{1, 5, 9},
				{10, 11, 13},
				{12, 13, 15}
			};
			int k = 8;
			
			cout << "矩阵:" << endl;
			for (auto& row : matrix) {
				for (int num : row) cout << num << " ";
				cout << endl;
			}
			cout << "第 " << k << " 小的元素: " 
			<< kthSmallest(matrix, k) << endl;
			// 输出: 13
			
			matrix = {{1, 2}, {1, 3}};
			k = 2;
			cout << "\n矩阵:" << endl;
			for (auto& row : matrix) {
				for (int num : row) cout << num << " ";
				cout << endl;
			}
			cout << "第 " << k << " 小的元素: " 
			<< kthSmallest(matrix, k) << endl;
			// 输出: 1
			
			return 0;
		}
	\end{lstlisting}
	
	\section{二分答案——一种重要的解题思想}
	
	\subsection{什么时候使用二分答案？}
	
	\begin{bluebox}
		\textbf{二分答案的典型场景}
		
		当一个问题满足以下条件时，可以尝试二分答案：
		\begin{enumerate}
			\item \textbf{答案具有单调性}：答案越大（或越小），越容易满足条件
			\item \textbf{可以快速验证}：对于给定的答案，能在较短时间内判断是否可行
			\item \textbf{求最值}：通常求"最大值的最小值"或"最小值的最大值"
		\end{enumerate}
	\end{bluebox}
	
	\subsection{二分答案的通用框架}
	
	\begin{lstlisting}[caption={二分答案的通用模板}]
		// 假设我们要找"满足条件的最小值"
		int binaryAnswer(long long left, long long right) {
			while (left < right) {
				long long mid = left + (right - left) / 2;  // 左边界模板
				
				if (check(mid)) {
					// mid满足条件，尝试更小的值
					right = mid;
				} else {
					// mid不满足条件，需要更大的值
					left = mid + 1;
				}
			}
			return left;  // 最小的满足条件的答案
		}
		
		// check函数：判断某个答案是否可行
		bool check(long long answer) {
			// 根据问题要求，判断以answer作为答案是否满足条件
			// 返回true表示满足，false表示不满足
		}
	\end{lstlisting}
	
	\section{总结：二分查找的核心要点}
	
	\begin{greenbox}
		\textbf{必须掌握的关键点}
		
		\begin{enumerate}
			\item \textbf{理解二段性}：二分查找的核心不是"有序"，而是能找到一种属性将数据分成两段
			\item \textbf{模板选择}：
			\begin{itemize}
				\item 找左边界（第一个满足条件的）：$mid = (l+r)/2, \; r = mid, \; l = mid+1$
				\item 找右边界（最后一个满足条件的）：$mid = (l+r+1)/2, \; l = mid, \; r = mid-1$
			\end{itemize}
			\item \textbf{防止死循环}：使用右边界模板时，$mid$计算必须$+1$
			\item \textbf{防止溢出}：使用 $mid = l + (r-l)/2$ 而不是 $(l+r)/2$
			\item \textbf{二分答案}：很多最优解问题可以通过二分答案转化为判定问题
			\item \textbf{灵活运用}：二分不限于数组，可以二分下标、二分值域、二分时间等
		\end{enumerate}
	\end{greenbox}
	
	\section{练习题推荐}
	
	\begin{enumerate}[leftmargin=*]
		\item \textbf{基础训练}：
		\begin{itemize}
			\item 查找插入位置（LeetCode 35）
			\item 找出有序数组中的单一元素（LeetCode 540）
			\item 寻找比目标字母大的最小字母（LeetCode 744）
		\end{itemize}
		\item \textbf{变形训练}：
		\begin{itemize}
			\item 搜索旋转排序数组 II（LeetCode 81）
			\item 寻找旋转排序数组中的最小值（LeetCode 153）
			\item 找出峰值（LeetCode 162）
		\end{itemize}
		\item \textbf{二分答案}：
		\begin{itemize}
			\item 吃香蕉的速度（LeetCode 875）
			\item 分割数组的最大值（LeetCode 410）
			\item 制作m天所需的花束（LeetCode 1482）
		\end{itemize}
		\item \textbf{困难挑战}：
		\begin{itemize}
			\item 寻找两个正序数组的中位数（LeetCode 4）
			\item 找出第k小的距离对（LeetCode 719）
			\item 有序矩阵中第K小的元素（LeetCode 378）
			\item 最小化最大值（LeetCode 410）
			\item 青蛙过河（LeetCode 403）
		\end{itemize}
	\end{enumerate}
	
\end{document}